\section{The $\alpha$-Dependent Lift}
\label{sec:lift}

We construct the right inverse of $\psia$ on its restriction to
$\polyringat{192}$. The opening protocol will lift each
$\GFt$-element appearing in the tensor-decomposed evaluation point
$q_0, q_1$ through this map; the only $\GFt$-elements that need the
lift are those drawn fresh from the transcript inside Step~7.

\subsection{Why the bit-pattern lift fails}

\begin{definition}[Bit-pattern lift]
\label{def:bitlift}
For $g \in \GFt$, the \emph{bit-pattern lift} of $g$ is its canonical
degree-$<192$ representative
\[
  g_{\mathrm{rep}}(X) \;=\; \sum_{i=0}^{191} g_i \cdot X^i \;\in\; \polyringat{192},
\]
where $g_i$ is the $i$-th bit of the binary expansion of $g$
(equivalently, the coefficient of $X^i$ in any polynomial in the
equivalence class of $g$ modulo $P$).
\end{definition}

\begin{remark}
\label{rem:bit-pattern-failure}
$\psia(g_{\mathrm{rep}}) = \sum_i g_i \cdot \alpha^i \in \GFt$. This
equals $g$ if and only if $\alpha$ is the equivalence class of the
formal $X$ in $\F_2[X]/(P)$ — that is, the field's quotient generator.
For a transcript-fresh $\alpha$ this happens with probability
$2^{-192}$.
\end{remark}

For a generic transcript-fresh $\alpha$ the bit-pattern lift therefore
does \emph{not} satisfy the round-trip $\psia \circ (\,\cdot\,)_{\mathrm{rep}}
= \mathrm{id}_{\GFt}$. The construction's verifier checks
$\psia(a') \stackrel{?}{=} a$, which collapses to the round-trip identity
under bit-pattern lift and the substitutions
$a = q_1^\top \cdot \psia(M_w) \cdot q_2$,
$a' = q_1'^\top \cdot M_w \cdot q_2'$ — so the bit-pattern lift is
unsound here.

\subsection{The polynomial basis and its inverse}

We replace the bit-pattern lift with the unique $\F_2$-linear
\emph{inverse-of-evaluation} lift. The key observation is that
$\GFt$ is a vector space of dimension $192$ over $\F_2$, and the
family
\[
  \mathcal{B}_\alpha \;:=\; \big\{1,\,\alpha,\,\alpha^2,\,\ldots,\,\alpha^{191}\big\}
\]
generates $\GFt$ as an $\F_2$-vector space for an overwhelming
fraction of $\alpha$.

\begin{lemma}[Generic invertibility]
\label{lem:generic-basis}
For $\alpha \in \GFt$ uniform, $\mathcal{B}_\alpha$ is an
$\F_2$-basis of $\GFt$ with probability at least
$1 - 2^{-95}$.
\end{lemma}

\begin{proof}
$\mathcal{B}_\alpha$ fails to be a basis iff $\alpha$ lies in a
proper subfield of $\GFt$. The proper subfields are
$\F_{2^d}$ for $d \mid 192, d < 192$. The largest is $\F_{2^{96}}$
(when $96 \mid 192$); its order is $2^{96}$. Including all proper
subfields and using the inclusion-exclusion bound
$\big|\bigcup_{d \mid 192, d < 192} \F_{2^d}\big| \le 2^{96} + 2^{64}
+ \cdots \le 2 \cdot 2^{96}$, the probability of bad $\alpha$ is
$\le 2 \cdot 2^{96} / 2^{192} = 2^{-95}$.
\end{proof}

For $\alpha$ with $\mathcal{B}_\alpha$ a basis, the $\F_2$-linear map
$\Phi_\alpha : \F_2^{192} \to \GFt$, $c \mapsto \sum_{i} c_i \alpha^i$,
is a bijection. We identify $\F_2^{192}$ with $\polyringat{192}$ by
$c \leftrightarrow \sum_i c_i X^i$; under this identification,
$\Phi_\alpha = \psia$ restricted to $\polyringat{192}$.

\begin{definition}[Inverse lift]
\label{def:invlift}
Let $\lift := \psia^{-1}|_{\polyringat{192}} : \GFt \to \polyringat{192}$
denote the unique $\F_2$-linear inverse of $\Phi_\alpha$. Explicitly,
for $g \in \GFt$, $\lift(g)$ is the unique
$\sum_i c_i X^i \in \polyringat{192}$ with
$\sum_i c_i \alpha^i = g$ in $\GFt$.
\end{definition}

\begin{proposition}[Defining property]
\label{prop:lift-roundtrip}
For every $g \in \GFt$,
$\psia\big(\lift(g)\big) = g$. The lift is $\F_2$-linear:
$\lift(g_1 + g_2) = \lift(g_1) + \lift(g_2)$.
\end{proposition}

\begin{proof}
First identity by definition. Linearity follows from
$\psia$ being $\F_2$-linear and $\lift$ being its inverse on a
basis.
\end{proof}

\subsection{Computing the lift}

We give a concrete algorithm. Let $M_\alpha \in \F_2^{192 \times 192}$
denote the matrix whose $j$-th \emph{column} is the bit-vector
representation of $\alpha^j \in \GFt$, $j = 0, \ldots, 191$. Then
\[
  M_\alpha \cdot c \;=\; (\Phi_\alpha(c))_{\mathrm{rep}}
  \qquad \text{for any column vector } c \in \F_2^{192}.
\]
By \Cref{lem:generic-basis}, $M_\alpha$ is invertible over $\F_2$
with probability $\ge 1 - 2^{-95}$ over $\alpha$, and the lift is
\[
  \lift(g) \;=\; M_\alpha^{-1} \cdot g_{\mathrm{rep}} \;\in\; \F_2^{192}.
\]

\begin{construction}[Inverse-lift table]
\label{cons:invlift}
On input $\alpha \in \GFt$:
\begin{enumerate}[leftmargin=2em]
  \item Compute the powers $1, \alpha, \alpha^2, \ldots, \alpha^{191}$
    by $191$ multiplications in $\GFt$, each at most a $192\times 192$
    bit carryless multiply followed by reduction modulo $P$.
  \item Assemble $M_\alpha \in \F_2^{192 \times 192}$, storing each
    column as three $64$-bit words.
  \item Compute $M_\alpha^{-1}$ over $\F_2$ by Gauss--Jordan
    elimination on an augmented $[M_\alpha \mid I]$. If the algorithm
    fails (pivoting cannot proceed), $\alpha$ is in the negligible
    fraction of \Cref{lem:generic-basis}; abort and resample.
  \item Per-entry, $\lift(g) = M_\alpha^{-1} \cdot g_{\mathrm{rep}}$
    is computed as $192$ bitwise inner products of $192$-bit row
    vectors, each in $O(1)$ word operations.
\end{enumerate}
The table $M_\alpha^{-1}$ is built once per transcript $\alpha$ and
shared between prover and verifier.
\end{construction}

\paragraph{Cost.} Construction~\ref{cons:invlift} runs in:
\begin{itemize}[leftmargin=2em]
  \item $191$ $\GFt$-multiplications (powers of $\alpha$) — bounded
    by a constant per multiplication;
  \item $\Theta(192^3 / 64) \approx 10^5$ word-level operations for
    the Gauss--Jordan inversion;
  \item $\Theta(192)$ word XORs per per-entry lift, of which there are
    at most $O(n_{\mathrm{rows}} + n_{\mathrm{cols}})$ in the
    opening (the lengths of the $q_0, q_1$ tensors).
\end{itemize}

\subsection{Soundness consequence}

\begin{corollary}[Lift commutes with $\psia$]
\label{cor:lift-commutes}
For any $u_1, \ldots, u_k \in \GFt$ and any polynomial-multiplication
expression $p(\cdot) \in \polyring[Y_1, \ldots, Y_k]$,
\[
  \psia\!\big(p\big(\lift(u_1), \ldots, \lift(u_k)\big)\big)
  \;=\;
  p(u_1, \ldots, u_k) \;\in\; \GFt,
\]
where the inner $p(\lift(u_1), \ldots, \lift(u_k))$ is evaluated in
$\polyring$ (no reduction) and the outer $\psia$ projects.
\end{corollary}

\begin{proof}
$\psia$ is an $\F_2$-algebra homomorphism, so it commutes with sums
and products in $\polyring$. By \Cref{prop:lift-roundtrip},
$\psia(\lift(u_i)) = u_i$ for every $i$. Therefore
$\psia(p(\lift(u_1), \ldots, \lift(u_k))) =
p(\psia(\lift(u_1)), \ldots, \psia(\lift(u_k))) =
p(u_1, \ldots, u_k)$.
\end{proof}

\Cref{cor:lift-commutes} is the operative form: \emph{any
polynomial expression in $\GFt$-elements equals the $\psia$-projection
of the same expression evaluated over $\polyring$ on the lifted
inputs.} The construction's verifier exploits this in
\Cref{sec:construction} to discharge the MLE-evaluation claim
$a = q_1^\top \cdot \psia(M_w) \cdot q_2$ by reducing it to the
identity $\psia(a') = a$ on a single lifted polynomial $a' \in
\polyring$ which the proximity check binds to the commitment.

\paragraph{Witness cells stay in $\cellring$.} A subtle but crucial
point: the witness cells $M_{w, b}$ in the matrix identity
$a = q_0^\top \cdot \psia(M_w) \cdot q_1$ are already polynomials in
$\polyringat{\Dval} \subset \polyring$; their $\psia$-image is
$\psia(M_{w,b}) = M_{w,b}(\alpha) \in \GFt$, which is the
\emph{evaluation} of the polynomial $M_{w,b}$ at $\alpha$ — exactly
what $\psia$ does on polynomial input. The lift therefore acts only on
the $q_0, q_1$ entries (the transcript-fresh $\GFt$-randomness); the
witness side carries through unchanged.
